% Experiments section -- filtered-proxy soft weighting for Asymmetric Duo TTA.
% Meant to be \input{} into the main report; assumes booktabs is loaded in
% the preamble (\usepackage{booktabs}). Replace the illustrative results in
% Table~\ref{tab:duo-results} with the full per-corruption sweep once it is
% run -- the numbers here come from single-seed exploratory runs on one duo
% and are meant to motivate the design choice, not to stand in for it.

\section{Experiments}
\label{sec:experiments}

\subsection{From frozen screening to adapting duos}
\label{sec:experiments:shift}

Our earlier experiments held both models frozen (\texttt{no\_adapt}) and used
this static setting purely as a cheap screen: for a fixed Asymmetric Duo, we
swept every combination of reliability proxy (\texttt{nuclear\_norm},
\texttt{atc}, \texttt{prototype}, \texttt{ac\_mc}, \texttt{cot}), Section-3
calibration map (\texttt{identity}, \texttt{linear}, \texttt{platt},
\texttt{beta}, \texttt{isotonic}), and proxy batch size $b_t$, and ranked the
resulting configurations by balanced selection accuracy
(\texttt{bal\_sel\_acc}) against the true per-chunk winner. This told us
\emph{which raw signal tracks reliability at all}, independent of whether
either model was actually adapting.

The project's focus has since shifted to the setting the method is actually
built for: both models undergoing live test-time adaptation (TENT, entropy
minimization) while the same filtered-proxy gate continuously re-weights
their combination. This changes the problem in two ways a frozen sweep
cannot capture. First, the input distribution to the proxy is no longer
fixed within a corruption: as each model's BatchNorm/LayerNorm statistics
adapt to the shift, the raw proxy score's own scale and dynamic range drift
over the course of the stream, so a proxy or calibration choice that looked
best on frozen logits is not guaranteed to remain best once the models it is
scoring are themselves changing under it. Second, the gate weight now feeds
back into what the models learn: whichever model the gate favors more
strongly shapes the entropy-minimization signal driving adaptation, a
coupling that has no analogue in the frozen, no-feedback setting. Every
result reported below is therefore from a live-adapting run
(\texttt{scripts/plot\_run\_diagnostics.py}, \texttt{mode=both\_indep} or
\texttt{both\_duo}), not the earlier frozen sweep.

\subsection{Current recipe: nuclear-norm proxy, identity calibration, EMA filtering}
\label{sec:experiments:recipe}

The configuration we have converged on for the adapting setting is:

\begin{itemize}
    \item \textbf{Proxy (Section~2):} \texttt{nuclear\_norm} -- the nuclear
    norm of the batch's softmax matrix. It is stateless (no source-data
    fitting pass is required, unlike \texttt{atc}/\texttt{prototype}/%
    \texttt{cot}), which matters more under adaptation than it did in the
    frozen sweep: there is no held-out clean pass to re-fit mid-stream, so a
    proxy that needs one is implicitly assuming its source-fitted state
    (e.g. \texttt{atc}'s threshold, \texttt{prototype}'s class means) stays
    valid as the models drift -- an assumption \texttt{nuclear\_norm} simply
    does not make.
    \item \textbf{Calibration (Section~3):} \texttt{identity} -- i.e., no
    fitted calibration map at all; the gate reads the raw proxy score
    directly. Section-3 calibration was introduced specifically to correct a
    proxy that is \emph{structurally} biased toward whichever model is
    stronger on average (see the \texttt{sel\_acc} vs.\ \texttt{bal\_sel\_acc}
    gotcha from the frozen-sweep phase). For \texttt{nuclear\_norm} on the
    duos we have adapted so far, this bias did not reappear badly enough to
    justify a fitted map: the identity baseline performed comparably to the
    learned calibration methods while removing an entire fitting stage (and
    its own held-out-corruption dependency) from the pipeline. We therefore
    drop calibration-map selection as an active axis of experimentation for
    now, rather than continuing to carry \texttt{linear}/\texttt{platt}/%
    \texttt{beta}/\texttt{isotonic} as live alternatives.
    \item \textbf{Temporal filter (Section~4):} exponential moving average
    (\texttt{ema}), $\alpha = 0.5$. This is a comparatively fast-reacting
    setting (half of the filter's state is replaced by each new
    observation), chosen because under live adaptation the model's own
    reliability can shift within a single corruption (e.g.\ as BN/LN
    statistics settle), and an overly slow filter lags exactly the kind of
    change the gate needs to react to.
    \item \textbf{Gate sharpness:} $\beta = 4.0$, fixed rather than fit via
    \texttt{fit\_beta} for the results below.
    \item \textbf{Proxy batch size:} $b_t = 128$, decoupled from the
    adaptation batch size (64 in our primary duo config) -- the gate
    recomputes every 128 samples rather than every adaptation step.
    \item \textbf{Base temperature prior:} a frozen \texttt{JointFixedTS}
    fit on clean validation data (\texttt{checkpoints/fixed\_ts/default})
    supplies $T_l, T_s$ for the Section-5 combination; the proxy sets only
    the mixing weight between the two temperature-scaled models, not their
    base scale.
\end{itemize}

\subsection{Illustrative results}
\label{sec:experiments:results}

Table~\ref{tab:duo-results} shows two representative corruptions from a
single-seed, live-adapting run of the ViT-B/16 (large) + ResNet-50 (small)
duo (\texttt{both\_indep}, i.e.\ each model adapts on its own logits while
the same nuclear-norm/identity/EMA gate combines them for the reported duo
accuracy), against two fixed baselines (\texttt{fixed\_ts}, \texttt{coca\_ts})
and an alternate filtered-proxy configuration
($b_t = 512$, $\beta$ fit via \texttt{fit\_beta} rather than fixed).

\begin{table}[h]
    \centering
    \caption{Per-corruption average accuracy, ViT-B/16 + ResNet-50 duo,
    severity 5, single seed. ``Ours'' is nuclear\_norm + identity + EMA,
    $b_t=128$, $\beta=4.0$ (fixed). ``Ours ($b_t{=}512$, fit $\beta$)'' is
    the same proxy/calibration/filter with a larger, less reactive proxy
    batch and a dev-corruption-tuned $\beta$.}
    \label{tab:duo-results}
    \begin{tabular}{lcccccc}
        \toprule
        Corruption & Large & Small & Ours & Fixed-TS & COCA-TS & Ours ($b_t{=}512$, fit $\beta$) \\
        \midrule
        brightness/s5        & 0.726 & 0.715 & \textbf{0.739} & --    & --    & -- \\
        elastic\_transform/s5 & 0.231 & 0.527 & \textbf{0.507} & 0.360 & 0.321 & 0.370 \\
        \bottomrule
    \end{tabular}
\end{table}

Two things stand out even from this small sample. First, the benefit of the
filtered-proxy gate is highly corruption-dependent: on
\texttt{brightness/s5}, where the two adapting models end up close in
accuracy (0.726 vs.\ 0.715), the gate's advantage over either input model is
marginal (0.739). On \texttt{elastic\_transform/s5}, where the large model
collapses under adaptation (0.231) while the small model does not
(0.527), the gate tracks the small model's reliability and recovers most of
its accuracy (0.507) -- exactly the regime the method is meant for, and
consistent with \texttt{screen\_duo\_candidates.py}'s notion of ``oracle
$\Delta$acc'': the achievable gain from a good per-chunk selector is largest
precisely when the two models disagree about who is reliable. Second, on
that same corruption, our fixed, small-$b_t$ configuration
(0.507) clearly outperforms both fixed baselines (0.360, 0.321) \emph{and}
the larger-$b_t$/fit-$\beta$ variant (0.370) -- i.e.\ reacting faster with a
fixed $\beta$ beat reacting slower with a supposedly better-tuned one on
this stream, which is itself a reason to treat $b_t$ and $\beta$ as tightly
coupled rather than independently optimal choices (see
Section~\ref{sec:experiments:future}).

These numbers are from one duo, one seed, and two corruptions, and are
reported to motivate the design choice in
Section~\ref{sec:experiments:recipe}, not as a full evaluation -- the
comprehensive per-corruption, multi-seed comparison across
\texttt{EVAL.CORRUPTIONS} is future work (Section~\ref{sec:experiments:future}).

\subsection{Future experiments}
\label{sec:experiments:future}

\begin{enumerate}
    \item \textbf{Kalman vs.\ EMA filtering.} The Kalman filter
    (\texttt{filter\_kind=kalman}) derives its gain from an explicit
    process/observation noise model ($q$, $r$) rather than a fixed EMA
    $\alpha$; in steady state it reduces to an EMA with a constant gain, but
    its transient behavior (how it responds right after a corruption
    boundary, before enough observations have accumulated) differs. Worth
    comparing directly against EMA under live adaptation, where the
    transient right after each corruption switch -- exactly when BN/LN
    statistics are also most in flux -- is arguably the most important part
    of the stream for the gate to get right.
    \item \textbf{Re-sweep $b_t$ under live adaptation.} The original
    proxy-batch-size sweep (\texttt{sweep\_proxies.py}) was run frozen. Given
    the elastic\_transform/s5 result above, where a small, fast-reacting
    $b_t$ outperformed a larger, fit-$\beta$ one, the reactivity/noise
    trade-off should be re-characterized specifically in the adapting
    regime, where the ``ground truth'' the gate is chasing (each model's
    instantaneous reliability) is itself moving as the models adapt --
    unlike the frozen setting, where it is static within a corruption.
    \item \textbf{Cross-duo confirmation.} All results so far are on ViT-B/16
    + ResNet-50. Prior sweeps found that the best proxy/calibration
    combination does \emph{not} transfer across duos (e.g.\ opposite bias
    directions for the same proxies between ResNet+ViT and
    ResNet+EfficientNet); the nuclear\_norm+identity+EMA recipe should be
    re-validated, not assumed, on the other five configured duos
    (ResNet-50+EfficientNet-B0, ConvNeXt+EfficientNet-B4, ConvNeXt+DeiT,
    ConvNeXt+Swin, ResNet-50+DeiT) before treating it as a general default.
    \item \textbf{Joint $\beta$/$b_t$ tuning under adaptation.} Rather than
    fixing $\beta=4.0$ and separately choosing $b_t$, grid-search both
    jointly (\texttt{fit\_beta} extended to sweep $b_t$ too) against
    held-out dev-corruption NLL computed \emph{with adaptation running}, not
    on frozen dev logits as \texttt{fit\_beta} currently does.
    \item \textbf{Headroom under adaptation.} Compare the current recipe
    against the two available upper bounds --
    \texttt{proxy\_kind=oracle} (best achievable through the sigmoid gate
    form) and \texttt{calibration\_mode=optimal\_w\_oracle}
    (best achievable $w_l$ with no functional-form constraint at all) -- now
    live-adapting, to quantify how much of the achievable gain the
    label-free recipe is actually capturing once adaptation feedback is in
    the loop, rather than relying on the frozen-setting headroom estimates.
    \item \textbf{Adaptation mode sensitivity.} Compare \texttt{both\_duo}
    (both models adapt on the joint calibrated signal) against
    \texttt{both\_indep} (each adapts independently, as used in
    Section~\ref{sec:experiments:results}) under the same recipe: adapting
    on the gate's own output couples the two models' training signals in a
    way independent adaptation does not, and it is not yet established
    which mode the current recipe is best tuned for.
    \item \textbf{Held-out generalization check.} Re-run the recipe under the
    \texttt{\_heldout} config variants (evaluating on the corruptions
    normally reserved for calibrator/dev fitting) to confirm the choice of
    nuclear\_norm+identity+EMA was not implicitly overfit to
    \texttt{EVAL.CORRUPTIONS} during exploration.
    \item \textbf{Data-quality fix before further EfficientNet runs.} An
    audit of the duo configs found every EfficientNet-B0/B4 duo config sets
    that side's TENT norm type to \texttt{LN}; torchvision's EfficientNet is
    pure \texttt{BatchNorm2d}, so this may mean \texttt{collect\_params} has
    been silently adapting zero parameters for EfficientNet in those runs.
    This should be corrected and those duos re-run before drawing
    cross-duo conclusions involving EfficientNet.
\end{enumerate}
